\chapter{Worked Example: From a Question to a Reviewed Report}\label{ch:example}

This chapter walks one small task from question to figure to report to review, using the
synthetic counting example of the repository \cite[Two renderers, one evidence
  contract]{example:custom-plotting}. It uses no measured data, so the task is tiny on
purpose: what matters is the order of the steps and the evidence each one leaves. The
same chapter is also how this handbook was built: it follows the report-writing
procedure, and its own review is the last step below.

\section{State the question and the contract}

\begin{headline}
  Write the stage record before drawing anything (Chapter~\ref{ch:code}).
\end{headline}

\textbf{Question.} How many of six synthetic entries fall into each of three fixed bins, and
does the plotted histogram show the same numbers as an independent hand count?

\begin{table}[H]
  \small
  \begin{tabularx}{\textwidth}{L{0.27\textwidth}Y}
    \toprule Field & Record before editing \\
    \midrule
    Authoritative source & The six literal samples in the recipe and the three bin edges
    $0,1,2,3$. \\
    \rowcolor{ctpSurface0!30}
    Named stages & Declare samples and edges; count per bin; compare with the hand count;
    draw; record the printed values. \\
    Oracle & A hand count written down independently of the code: $[2,3,1]$. \\
    \rowcolor{ctpSurface0!30}
    Acceptance evidence & Counts equal the oracle; the rendered page carries the axis
    labels and the printed bin values; the render gate has no blocking finding. \\
    Non-goals & No measured data, uncertainty, fit or physics selection. \\
    \bottomrule
  \end{tabularx}
  \caption{The stage record of the example, written before the recipe is run
    \cite[Establish the scientific contract]{skill:coding}.}
  \label{tab:example-stage}
\end{table}

The plotted quantity is an equation first. With samples $x_j$ and edges
$e_0<e_1<e_2<e_3$, the content of bin $i$ is
\begin{equation}
  N_i = \sum_{j=1}^{6} \mathbf{1}\!\left[e_{i-1}\le x_j<e_i\right],
  \qquad i=1,2,3,
  \label{eq:counts}
\end{equation}
where $\mathbf{1}[\cdot]$ is 1 when sample $x_j$ lies in bin $i$ and 0 otherwise. The figure
shows $N_i$ as a function of the coordinate, with $N_i$ on its vertical axis \cite[Design
  and iterate at two scales]{skill:plot-production}.

\section{Run the recipe and keep its evidence}

\begin{headline}
  The recipe checks its own counts against the oracle before it draws, and leaves its
  evidence next to the figure.
\end{headline}

From the repository root, with a Python environment that has NumPy and Matplotlib:
\begin{lstlisting}
python examples/custom-plotting/render.py --renderer matplotlib \
  --out docs/guide/build/custom
\end{lstlisting}

\begin{figure}[H]
  \centering
  \begin{tikzpicture}[node distance=6mm and 9mm]
    \node[work,text width=24mm] (in) {Samples and edges};
    \node[work,text width=24mm,right=of in] (n) {Count per bin};
    \node[check,text width=32mm,right=of n] (o) {Equal to $[2,3,1]$, sum equal to the
      number of samples?};
    \node[result,text width=24mm,right=of o] (draw) {Draw and write evidence};
    \node[box,draw=ctpText!60,text width=24mm,below=7mm of o] (stop) {Stop};
    \draw[flow] (in)--(n);
    \draw[flow] (n)--(o);
    \draw[flow] (o)--node[above,lbl]{yes}(draw);
    \draw[flow] (o)--node[right,lbl]{no}(stop);
  \end{tikzpicture}
  \caption{What the recipe does.}
  \label{fig:example-recipe}
\end{figure}

The evidence written next to the figure:
\begin{itemize}
  \item the source samples and the formatted values;
  \item the renderer identity;
  \item an inventory of the rendered page;
  \item an applicability record for the render gate.
\end{itemize}
The recipe refuses to reuse an output directory, so old evidence cannot be mixed with new.
Figure~\ref{fig:example} is its output.

\begin{figure}[H]
  \centering
  \includegraphics[width=0.8\textwidth]{build/custom/example-counts.pdf}
  \caption{Synthetic example: entries $N_i$ per bin, as in
    Equation~\ref{eq:counts}. Read the three printed bin values, 2, 3 and 1, and the total
    of 6 at the top right: they agree with the independent hand count. There are no
    measured data.}
  \label{fig:example}
\end{figure}

\evidence\ The oracle and the computed counts agree: $N=(2,3,1)$, total $6$. This supports
only the declared coverage, which is these six samples and three bins
\cite[Numerical evidence]{skill:implementation}.

\section{Write the report around the figure}

\begin{headline}
  A report for this question can be one page, and its claim is limited to the synthetic
  input.
\end{headline}

The one-page report \cite[Source traceability and corroboration; Design the
  argument]{skill:report-production}:
\begin{itemize}
  \item opens with the question, the answer and its status;
  \item introduces Figure~\ref{fig:example} before interpreting it, states what is held fixed
        (samples and edges) and which feature answers the question (the three printed bin
        values);
  \item cites the source that owns each claim: the recipe's evidence for the counts, the stage
        record for the contract, a workflow document only for why the checks come in this order.
\end{itemize}

\gap\ A report that said ``the method works'' would claim more than a synthetic check can
show. The claim that survives is narrower: the recipe's counts match a hand count for these
inputs, and the page carries the values it should.

\section{Review in separate tracks}

\begin{headline}
  Each review track of Chapter~\ref{ch:reports} answers one question. For this handbook
  they run as follows, from the repository root.
\end{headline}

\begin{description}
  \item[Source] \code{skills/report-writing/scripts/check_report_sources.py} checks the chapter
        sources for duplicate bibliography keys, \texttt{\textbackslash addbibresource} targets
        outside the shared library, hand-written bibliography blocks and formatting drift. A
        reader verifies pinpoints by following each citation to its entry.
  \item[Build and bibliography] \code{latexmk} runs \code{pdflatex} and \code{biber} until
        references settle; the \code{.log} and \code{.blg} files are read for unresolved keys or
        references.
  \item[Rendered pages] \code{skills/render-review/scripts/render_gate.py gate} binds the PDF,
        its log and an applicability contract by SHA-256 and writes \code{gate.json} and a review
        packet. The contract records, for each optional check, why it does not apply here. A
        \code{fail} blocks visual review until it is fixed or waived \cite[Command;
          Rules]{skill:render-workflow}.
  \item[Reading tasks] After the gate, a reviewer receives selected page images, never the PDF,
        with the tasks to perform: read the axis labels in Figure~\ref{fig:example}, find the bin
        with three entries, and say what the page claims \cite[Figure and page
          tasks]{skill:visual-review}.
\end{description}

\review\ The gate's status is mechanical evidence about layout, fonts, build errors and
printed values. It is not a verdict that the page reads well; the reading tasks supply that.
The commands and the location of the gate output are in the \code{README.md} beside this
handbook's sources.
